%% 第4巻　形と向き　本文（第13回編集委員会）
%% 章は \bkinputchap{ch4-NN} で読む（まだない章は、組むときに飛ばして「BKPAGE missing」を出す）
%% 第4巻は巻の色が黄緑なので、部扉の絵の線は白でなく bkOn（濃い紺）でかく
\providecommand\bkinputchap[1]{\IfFileExists{#1.tex}{\input{#1}}{\clearpage\typeout{BKPAGE missing #1}}}
\bkinputchap{ch4-00}

%% ======================== 第1部 図形 ========================
\bkPartOpener{1}{第4巻\hspace{0.8em}第1部}{図形}%
  {172/{三角形の角をたすと、なぜ必ず180°？},
   160/{影の長さから、木の高さがわかる？},
   148/{ななめの距離を、たてと横から計算できる？}}%
  {% 直角三角形と正方形（三平方の定理。濃い紺の線）
   \begin{scope}[bkOn, line width=1.2pt, line join=round, shift={(72,57)}, scale=0.55]% 第13回：帯の中、問いの字から離した右下へ
     \draw (130,118) -- (154,118) -- (154,136) -- cycle;
     \draw[line width=0.9pt] (130,118) -- (154,118) -- (154,94) -- (130,94) -- cycle;
     \draw[line width=0.9pt] (154,118) -- (172,118) -- (172,136) -- (154,136) -- cycle;
     \draw[line width=0.9pt] (130,118) -- (154,136) -- (136,160) -- (112,142) -- cycle;
     \draw[line width=0.8pt] (151,118) -- (151,121) -- (154,121);
   \end{scope}}%
  {形のきまりを、図をかいて確かめ、ことばで説明できるようにするよ。にた形の比と三平方の定理を使うと、行けない場所の長さも計算できるよ。}
\bkPartContents{1}{%% 自動生成（gen.py）
\bktocchap{4-1}{きまりから考えを積み上げる：定義・公理・定理}{三角形の角をたすと、なぜ必ず180°？}{6}{2}{3}{26}{C}
\bktocchap{4-2}{角度と円、ラジアン}{角度を「長さ」ではかる方法}{4}{1}{2}{56}{C}
\bktocchap{4-3}{にた形と比：相似}{影の長さから、木の高さがわかる}{6}{2}{3}{98}{C}
\bktocchap{4-4}{三平方の定理}{ななめの距離を、たてと横から計算する}{6}{3}{3}{128}{C}
\bktocchap{4-5}{面積と体積：細かく分けてたす}{円を切って並べかえると、長方形になる？}{2}{1}{2}{160}{C}
}%
{角度・面積・体積と、にた形の比、三平方の定理を使って、ものの大きさや距離を計算できるようになります。「細かく分けてたす」面積の考え方は、第5巻の積分の入口です。}
\bkpartbegin{1}{図形}
\bkinputchap{ch4-01}
\bkinputchap{ch4-02}
\bkinputchap{ch4-03}
\bkinputchap{ch4-04}
\bkinputchap{ch4-05}

%% ======================== 第2部 三角比と三角関数 ========================
\bkPartOpener{2}{第4巻\hspace{0.8em}第2部}{三角比と三角関数}%
  {172/{坂の上の荷物にかかる力を、2つに分けられる？},
   160/{回る点の影は、なぜ行ったり来たりする？},
   148/{2つの音を重ねると、なぜウワンウワンと聞こえる？}}%
  {% 円と、回る点の影の波（濃い紺の線）
   \begin{scope}[bkOn, line width=1.2pt, line cap=round, line join=round]
     \draw (130,128) circle (14);
     \draw (116,128) -- (172,128);
     \draw[line width=0.9pt] (130,128) -- ({130+14*cos(50)},{128+14*sin(50)});
     \fill ({130+14*cos(50)},{128+14*sin(50)}) circle (2);
     \draw[line width=1.6pt, domain=0:360, samples=60] plot ({146+\x/360*26}, {128+14*sin(\x+50)});
   \end{scope}}%
  {直角三角形の辺の比から、角度と長さを結びつけるよ。くるくる回る点の影を追いかけると、波の形のグラフ（三角関数）が現れるよ。}
\bkPartContents{2}{%% 自動生成（gen.py）
\bktocchap{4-6}{坂道と三角比}{坂の上の荷物にかかる力を、2つに分ける}{4}{2}{3}{196}{C}
\bktocchap{4-7}{三角形を解く：正弦定理と余弦定理}{行けない場所までの距離をはかる}{4}{3}{3}{234}{C}
\bktocchap{4-8}{三角関数：円運動の影}{回る点の影は、行ったり来たりする}{4}{3}{4}{268}{C}
\bktocchap{4-9}{波をたし合わせる：加法定理とうなり}{2つの音を重ねると、ウワンウワンと聞こえる}{4}{4}{4}{306}{C}
}%
{坂道や力の向きを三角比で計算し、回転や波を三角関数で表せるようになります。三角関数は、第5巻の微分、第6巻の振動、第9巻のフーリエ解析の主役です。}
\bkpartbegin{2}{三角比と三角関数}
\bkinputchap{ch4-06}
\bkinputchap{ch4-07}
\bkinputchap{ch4-08}
\bkinputchap{ch4-09}

%% ======================== 第3部 ベクトル ========================
\bkPartOpener{3}{第4巻\hspace{0.8em}第3部}{ベクトル}%
  {172/{2人で引くと、荷物はどっちへ動く？},
   160/{進む向きと直角に押しても、仕事は0？},
   148/{部屋の中の1点を、3つの数で表せる？}}%
  {% 2本の矢印とその和（濃い紺の線）
   \begin{scope}[bkOn, line width=1.4pt, line cap=round, line join=round]
     \draw[-{Stealth[length=3mm]}] (120,110) -- (150,118);
     \draw[-{Stealth[length=3mm]}] (120,110) -- (130,140);
     \draw[line width=0.8pt, dash pattern=on 2pt off 2pt] (150,118) -- (160,148) -- (130,140);
     \draw[-{Stealth[length=3.4mm]}, line width=2pt] (120,110) -- (160,148);
   \end{scope}}%
  {向きと大きさをもつ量（ベクトル）を、矢印と数の組で表すよ。たし算・内積・外積を使うと、力・速さ・回す力を計算できるよ。}
\bkPartContents{3}{%% 自動生成（gen.py）
\bktocchap{4-10}{向きと大きさをもつ量：ベクトル}{2人で引くと、荷物はどっちへ動く？}{6}{2}{2}{340}{C}
\bktocchap{4-11}{ベクトルを成分で計算しよう}{川に流されながら渡る船は、どこに着く？}{6}{2}{3}{374}{C}
\bktocchap{4-12}{内積：どれだけ同じ向きか}{進む向きと直角に押しても、仕事は0}{6}{3}{4}{408}{C}
\bktocchap{4-13}{空間の座標と3次元のベクトル}{部屋の中の1点を、3つの数で表す}{6}{4}{4}{434}{C}
\bktocchap{4-14}{外積：回す力と右ねじ}{ドアのノブは、なぜちょうつがいから遠い？}{6}{4}{5}{468}{C}
\bktocchap{4-15}{いろいろな座標：極座標・円柱座標・球座標}{地球の上の場所は、緯度と経度で表す}{6}{4}{5}{500}{C}
}%
{力や速さを成分に分けて計算し、仕事（内積）と回す力（外積）を求められるようになります。第7巻の行列、第8巻のベクトル解析は、ここから始まります。}
\bkpartbegin{3}{ベクトル}
\bkinputchap{ch4-10}
\bkinputchap{ch4-11}
\bkinputchap{ch4-12}
\bkinputchap{ch4-13}
\bkinputchap{ch4-14}
\bkinputchap{ch4-15}

%% ======================== 巻末 ========================
\bkinputchap{ch4-back}
